\section{Introduction}
\pagenumbering{arabic}

\noindent

Afin d'assurer le bon fonctionnement du réseau électrique, il est nécessaire d'avoir un équilibre entre la production et la consommation d'électricité. La consommation d'électricité est influencée par de nombreux facteurs, tels que la météo, les habitudes de consommation des ménages et des entreprises, et les événements spéciaux. La prévision de la consommation d'électricité est donc un enjeu crucial pour les gestionnaires de réseau: elle permet d'anticiper les besoins, de planifier la production et donc de de garantir la stabilité du système électrique. 

Afin de suivre la consommation électrique en temps réel, le Réseau de Transport d’Électricité RTE collecte les données de consommation à l'aide de télémesures transmises directement depuis ses installations \parencite{rte_eco2mix_2026}. Certains de ses sites ne sont pas équipé pour la télémesures, les données sont donc complétées par des estimations, remplacées par la suite par des données de comptage plus complètes. RTE propose également des prévisions de consommation à court terme, qui sont utilisées pour planifier la production et la distribution d'électricité. Ces prévisions sont basées sur des modèles statistiques et des algorithmes de machine learning, qui prennent en compte les données historiques de consommation, les conditions météorologiques et d'autres facteurs pertinents.

L'objectif de ce rapport est de proposer un algorithme de prédiction de la consommation électrique française à court terme, en utilisant les données de consommation collectées par RTE. Nous allons explorer différentes approches de modélisation, en utilisant des techniques d'analyse de séries temporelles et de machine learning, afin d'obtenir des prévisions précises et fiables. Nous allons également évaluer les performances de notre algorithme en le comparant à des benchmarks, des modèles statistiques et ainsi qu'à d'autres modèles de machine learning, et en analysant les erreurs de prévision pour identifier les points faibles et les opportunités d'amélioration. Nos prédictions s'appuieront sur les données historiques compilées par RTE, ainsi que sur les données météorologiques du réseau SYNOP \parencite{meteo_france_synop_2026}, et le calendrier des jours fériés et vacances scolaires en France.
